% !PRINTED 694
\weberpage{694}
\begin{center}
§. 191.

\gesperrt{Theorems from the theory of series.}
\end{center}

For the further applications of the results obtained hitherto, certain theorems from the theory of infinite series are required, and these shall first be derived here.

1. \gesperrt{If $a_1, a_2, a_3, \ldots, a_n\ldots$ is an unlimited system of real (positive, negative, or also vanishing) quantities, of which we assume that the sum}

\weq{1}{\sigma_n = a_1+a_2+\cdots+a_n}

\gesperrt{remains, for every arbitrary $n$, in absolute value below a finite limit $C$, then the series}

\weq{2}{S = \frac{a_1}{1^s}+\frac{a_2}{2^s}+\frac{a_3}{3^s}+\cdots}

\gesperrt{is, for every positive $s$, not only convergent, but also a continuous function of $s$.}

To prove this theorem, we decompose $S$ in the following way:
\[
S = S_m+R_m,
\]
in which
\[
S_m = \frac{a_1}{1^s}+\frac{a_2}{2^s}+\cdots+\frac{a_{m-1}}{(m-1)^s},
\]
\[
R_m = \frac{a_m}{m^s}+\frac{a_{m+1}}{(m+1)^s}+\frac{a_{m+2}}{(m+2)^s}+\cdots
\]
Now by (1):
\[
a_m = \sigma_m-\sigma_{m-1}, \quad a_{m+1}=\sigma_{m+1}-\sigma_m, \quad a_{m+2}=\sigma_{m+2}-\sigma_{m+1}, \ldots
\]
and therefore
\[
\begin{aligned}
R_m+\frac{\sigma_{m-1}}{m^s}
&= \sigma_m\left(\frac{1}{m^s}-\frac{1}{(m+1)^s}\right) \\
&\quad +\sigma_{m+1}\left(\frac{1}{(m+1)^s}-\frac{1}{(m+2)^s}\right)+\cdots
\end{aligned}
\]
Since now, by the assumption, $\sigma_{m-1}, \sigma_m, \sigma_{m+1}, \ldots$ are enclosed between finite limits $\pm C$, it follows from this that $R_m+\dfrac{\sigma_{m-1}}{m^s}$ is contained between the two limits
\[
\pm C\left(\frac{1}{m^s}-\frac{1}{(m+1)^s}+\frac{1}{(m+1)^s}-\frac{1}{(m+2)^s}+\cdots\right) = \pm\frac{C}{m^s}
\]
%CONT>>
